\documentclass[12pt, a4paper]{book}

\usepackage[T1]{fontenc}
\usepackage[utf8]{inputenc}
\usepackage[charter, cal=cmcal]{mathdesign}
\usepackage[scaled=0.92]{helvet}
\usepackage{geometry}
\geometry{a4paper, top=30mm, bottom=30mm, inner=32mm, outer=25mm, marginparwidth=18mm}

\usepackage{amsmath, amsthm}
\usepackage[version=4]{mhchem}
\usepackage{graphicx}
\usepackage{xcolor}
\usepackage{tikz}
\usetikzlibrary{arrows.meta, shapes.geometric, positioning, calc, backgrounds}
\usepackage{pgfplots}
\pgfplotsset{compat=1.18}
\usepackage{tcolorbox}
\tcbuselibrary{skins, breakable, theorems}
\usepackage{titlesec}
\usepackage{fancyhdr}
\usepackage{enumitem}
\usepackage{booktabs}
\usepackage{array}
\usepackage{parskip}
\usepackage{siunitx}
\usepackage[colorlinks=true, linkcolor=BrickRed,
            urlcolor=NavyBlue, citecolor=ForestGreen]{hyperref}

% ── Colours (Charcoal / Royal Blue) ──────────────────────────────────────────
\definecolor{JourneyColor}{HTML}{0F172A}
\definecolor{ModuleColor}{HTML}{2563EB}
\definecolor{PuzzleColor}{HTML}{D97706}
\definecolor{DiscoveryColor}{HTML}{3B82F6}
\definecolor{WarningColor}{HTML}{991B1B}
\definecolor{ProblemColor}{HTML}{6B21A8}
\definecolor{ReflectColor}{HTML}{0D9488}
\definecolor{HistoryColor}{HTML}{78350F}

% ── Boxes ─────────────────────────────────────────────────────────────────────
\newtcolorbox{puzzlebox}[1][]{enhanced,breakable,
  colback=PuzzleColor!6,colframe=PuzzleColor,fonttitle=\bfseries\sffamily,
  title={The Puzzle},attach boxed title to top left={yshift=-2mm,xshift=6mm},
  boxed title style={colback=PuzzleColor},borderline west={3pt}{0pt}{PuzzleColor},#1}

\newtcolorbox{discoverybox}[1][]{enhanced,breakable,
  colback=DiscoveryColor!8,colframe=DiscoveryColor,fonttitle=\bfseries\sffamily,
  title={Discovery},attach boxed title to top left={yshift=-2mm,xshift=6mm},
  boxed title style={colback=DiscoveryColor},borderline west={3pt}{0pt}{DiscoveryColor},#1}

\newtcolorbox{problembox}[1][]{enhanced,breakable,
  colback=ProblemColor!5,colframe=ProblemColor,fonttitle=\bfseries\sffamily,
  title={Problems},attach boxed title to top left={yshift=-2mm,xshift=6mm},
  boxed title style={colback=ProblemColor},borderline west={3pt}{0pt}{ProblemColor},#1}

\newtcolorbox{deepbox}[1][]{enhanced,breakable,
  colback=JourneyColor!5,colframe=JourneyColor,fonttitle=\bfseries\sffamily,
  title={Deep Problems},attach boxed title to top left={yshift=-2mm,xshift=6mm},
  boxed title style={colback=JourneyColor},borderline west={4pt}{0pt}{JourneyColor},#1}

\newtcolorbox{historybox}[1][]{enhanced,breakable,
  colback=HistoryColor!7,colframe=HistoryColor!60,fonttitle=\bfseries\sffamily\small,
  title={Historical Note},attach boxed title to top left={yshift=-2mm,xshift=6mm},
  boxed title style={colback=HistoryColor!80},#1}

\newtcolorbox{reflectionbox}[1][]{enhanced,breakable,
  colback=ReflectColor!6,colframe=ReflectColor,fonttitle=\bfseries\sffamily,
  title={Reflection},attach boxed title to top left={yshift=-2mm,xshift=6mm},
  boxed title style={colback=ReflectColor},borderline west={3pt}{0pt}{ReflectColor},#1}

\newtcolorbox{namedconcept}[2][]{enhanced,
  colback=ModuleColor!6,colframe=ModuleColor,fonttitle=\bfseries\sffamily\large,
  title={#2},attach boxed title to top left={yshift=-2mm,xshift=6mm},
  boxed title style={colback=ModuleColor},#1}

% ── Section formatting ────────────────────────────────────────────────────────
\titleformat{\chapter}[display]
  {\normalfont\sffamily\huge\bfseries\color{JourneyColor}}
  {\textcolor{JourneyColor!60}{\chaptertitlename\ \thechapter}}{10pt}{\Huge}
  [\vspace{2pt}\textcolor{JourneyColor}{\titlerule[1.5pt]}]
\titleformat{\section}
  {\normalfont\sffamily\Large\bfseries\color{ModuleColor}}
  {\thesection}{1em}{}[\textcolor{ModuleColor!40}{\titlerule[0.6pt]}]
\titleformat{\subsection}
  {\normalfont\sffamily\large\bfseries\color{JourneyColor!80}}{\thesubsection}{1em}{}
\titleformat{\subsubsection}
  {\normalfont\sffamily\normalsize\bfseries}{\thesubsubsection}{1em}{}

\pagestyle{fancy}
\fancyhf{}
\fancyhead[LE]{\small\sffamily\textcolor{JourneyColor}{\leftmark}}
\fancyhead[RO]{\small\sffamily\textcolor{ModuleColor}{\rightmark}}
\fancyfoot[C]{\small\thepage}
\renewcommand{\headrulewidth}{0.4pt}

\newcommand{\newterm}[1]{\textit{\textcolor{DiscoveryColor}{#1}}}
\newcommand{\important}[1]{\textcolor{ModuleColor}{\textbf{#1}}}

\begin{document}

\setcounter{chapter}{4}
\chapter[Journey 5 --- Algorithms, Data Structures, and Optimization]
        {Journey Five\\[8pt]
         \textcolor{ModuleColor}{\large Algorithms, Data Structures, and Optimization}}
\setcounter{section}{0}

\section{Module 5.1 --- Algorithmic Analysis and Divide-and-Conquer}
\label{sec:mod5.1}

\subsection*{Motivation}

We completed in Journey 4 the software-hardware boundary, studying how compilers translate syntactic structures into registers and machine instructions.
Now, we shift focus from the translation of code to the structural efficiency of the logic itself.
We begin our study of algorithmic analysis.

As the size of input data grows, the execution speed of a program is determined not by the speed of the compiler or hardware, but by the mathematical complexity class of the algorithm.
This module formalizes asymptotic notation (Big-O, $\Omega$, $\Theta$) to measure execution scaling.
We will study the **Divide-and-Conquer** design paradigm, formalize the Master Theorem to solve recurrences, and derive a fundamental mathematical proof: the absolute, comparison-based sorting lower bound of $\Omega(n \log n)$.

\subsection{The Sorting Complexity Boundary}

We investigate why divide-and-conquer reduces computation steps.

\begin{puzzlebox}
\textbf{The Comparison Limit.}

Suppose we want to sort an unsorted list of $n$ numbers.
* A basic sorting algorithm (like Bubble Sort or Insertion Sort) compares adjacent elements and swaps them, executing $O(n^2)$ comparisons.
* A divide-and-conquer algorithm (like Merge Sort) splits the list in half, recursively sorts each half, and merges the sorted halves in linear time, executing $O(n \log n)$ comparisons.
\begin{enumerate}[label=(\alph*), leftmargin=*]
  \item If $n = 10^6$ (one million elements):
        * Calculate the approximate number of comparisons executed by Insertion Sort ($n^2$).
        * Calculate the approximate number of comparisons executed by Merge Sort ($n \log_2 n$).
        If a CPU can execute $10^9$ comparisons per second, calculate the running time of both algorithms. Why does the complexity difference represent the boundary between usability and failure?
  \item We model any comparison-based sorting algorithm as a binary **Decision Tree**, where each node is a comparison ($x_i < x_j$), and each leaf is a unique sorted permutation of the input elements.
        For a list of $n$ elements, how many unique permutations (leaves) must the decision tree contain?
  \item The height of the decision tree represents the worst-case number of comparisons.
        Prove that a binary tree with $L$ leaves must have a height $h \ge \log_2 L$.
        Using Stirling's Approximation ($\ln n! \approx n \ln n - n$), show why the worst-case height of any comparison-based sorting decision tree must be at least $\Omega(n \log n)$.
        Does there exist any comparison-based sorting algorithm that can sort in $O(n)$ time?
\end{enumerate}
\end{puzzlebox}

\subsection{Asymptotic Complexity and the Master Theorem}

We formalize how to measure scaling and solve recursive recurrences.

\begin{namedconcept}{Asymptotic Definitions}
Let $f(n)$ and $g(n)$ be functions mapping positive integers to positive reals.
* **$f(n) = O(g(n))$ (Upper Bound):** There exist constants $c > 0$ and $n_0 > 0$ such that $f(n) \le c \cdot g(n)$ for all $n \ge n_0$.
* **$f(n) = \Omega(g(n))$ (Lower Bound):** There exist constants $c > 0$ and $n_0 > 0$ such that $f(n) \ge c \cdot g(n)$ for all $n \ge n_0$.
* **$f(n) = \Theta(g(n))$ (Tight Bound):** $f(n) = O(g(n))$ and $f(n) = \Omega(g(n))$.
\end{namedconcept}

\begin{explobox}
Many divide-and-conquer algorithms yield recurrence relations of the form:
\[
  T(n) \;=\; a \cdot T(n/b) + f(n)
\]
representing a problem of size $n$ divided into $a$ subproblems of size $n/b$, with $f(n)$ work done to split and merge.
We solve these recurrences using the **Master Theorem**:
1. If $f(n) = O(n^{\log_b a - \epsilon})$ for some $\epsilon > 0$, then $T(n) = \Theta(n^{\log_b a})$.
2. If $f(n) = \Theta(n^{\log_b a})$, then $T(n) = \Theta(n^{\log_b a} \log n)$.
3. If $f(n) = \Omega(n^{\log_b a + \epsilon})$ for some $\epsilon > 0$, and if $a \cdot f(n/b) \le c \cdot f(n)$ for some $c < 1$ and large $n$, then $T(n) = \Theta(f(n))$.
\begin{enumerate}[label=(\alph*), leftmargin=*]
  \item Solve the recurrence $T(n) = 4 T(n/2) + n$ using the Master Theorem. Specify which case applies.
  \item Solve the recurrence for Merge Sort: $T(n) = 2 T(n/2) + \Theta(n)$.
  \item Consider the recurrence $T(n) = 2 T(n/2) + n \log n$.
        Explain why the standard Master Theorem does not apply directly to this recurrence.
\end{enumerate}
\end{explobox}

\subsection{Pivot Selection and Randomized Quicksort}

We study the impact of pivot selection on quicksort complexity.

\begin{explobox}
Unlike Merge Sort, **Quicksort** partitions the array around a selected **pivot** element.
* **Best Case:** The pivot divides the array into two equal halves at each step $\implies T(n) = 2 T(n/2) + \Theta(n) \implies \Theta(n \log n)$.
* **Worst Case:** The pivot is always the smallest or largest element, dividing the array into $0$ and $n-1$ elements $\implies T(n) = T(n-1) + \Theta(n)$.
\begin{enumerate}[label=(\alph*), leftmargin=*]
  \item Solve the worst-case recurrence $T(n) = T(n-1) + c \cdot n$ using expansion. Show that it yields $\Theta(n^2)$.
  \item To prevent this worst-case behavior on sorted arrays, we use **Randomized Quicksort**, selecting the pivot uniformly at random from the subarray.
        Explain how selecting a random pivot ensures that we achieve $O(n \log n)$ average-case running time, regardless of the initial order of the input.
  \item Why is quicksort often faster in practice than merge sort, even though both have $O(n \log n)$ average time? (Think about cache locality and in-place memory usage).
\end{enumerate}
\end{explobox}

\newpage

% ─── Section: Discovery ──────────────────────────────────────────────────────
\section{Part 3 --- Discovery}

\begin{discoverybox}
\textbf{What have we built?}

\textbf{First.} Asymptotic notations (Big-O, $\Omega$, $\Theta$) mathematically define how algorithm execution times scale as inputs approach infinity.

\textbf{Second.} The Master Theorem provides a recipe for solving divide-and-conquer recurrences by comparing subproblem branching to split/merge costs.

\textbf{Third.} Comparison-based sorting has a mathematically provable worst-case lower bound of $\Omega(n \log n)$, derived by modeling sorting as a binary decision tree.

\textbf{Fourth.} Randomized pivot selection in Quicksort prevents $O(n^2)$ worst-case execution paths, guaranteeing $O(n \log n)$ average-case runtime.
\end{discoverybox}

\newpage

% ─── Section: Naming Concepts ────────────────────────────────────────────────
\section{Part 4 --- What This Is Called}

\begin{namedconcept}{Algorithmic parameters}
* **Asymptotic Analysis:** The study of how execution parameters (time, space) scale with input size.
* **Divide-and-Conquer:** An algorithmic paradigm that divides a problem into smaller subproblems, solves them recursively, and combines the results.
* **Master Theorem:** A theorem for solving asymptotic recurrence relations of divide-and-conquer algorithms.
* **Decision Tree:** A directed model mapping all possible sequences of comparisons executed by a sorting algorithm.
* **Randomized Algorithm:** An algorithm that makes random choices during execution to optimize average-case complexity.
\end{namedconcept}

\newpage

% ─── Section: Solved Exercises ───────────────────────────────────────────────
\section{Part 5 --- Solved Exercises}

\begin{problembox}
\begin{enumerate}[label=\textbf{E\arabic*.}]
  \item \textbf{Solving Recurrences via the Master Theorem.}
  
        Solve the recurrence relation $T(n) = 4 T(n/2) + n$ using the Master Theorem.
        
        \textbf{Solution:}
        The recurrence is:
        \[
          T(n) \;=\; a \cdot T(n/b) + f(n)
        \]
        where:
        * $a = 4$ (number of subproblems).
        * $b = 2$ (factor by which input size is divided).
        * $f(n) = n$ (work done outside recursive calls).
        
        1. **Compute the critical exponent $\log_b a$:**
           \[
             \log_b a \;=\; \log_2 4 \;=\; 2.
           \]
           The comparing function is $n^{\log_b a} = n^2$.
           
        2. **Compare $f(n)$ to $n^2$:**
           We compare $f(n) = n$ to $n^{\log_b a} = n^2$.
           Clearly, $f(n) = O(n^{2 - \epsilon})$ for $\epsilon = 1 > 0$.
           
        3. **Apply the Master Theorem:**
           This matches **Case 1** of the Master Theorem.
           The recursive branching dominates the complexity.
           Therefore:
           \[
             T(n) \;=\; \Theta\left( n^{\log_b a} \right) \;=\; \Theta\left( n^2 \right).
           \]
           Thus, the recurrence solves to $\Theta(n^2)$.
           
  \item \textbf{Stirling's Approximation and the Sorting Bound.}
  
        Prove that $\log_2(n!) = \Theta(n \log n)$ using Stirling's Approximation.
        
        \textbf{Solution:}
        Stirling's Approximation states that:
        \[
          n! \;\approx\; \sqrt{2 \pi n} \left( \frac{n}{e} \right)^n.
        \]
        Taking the natural logarithm of both sides:
        \[
          \ln(n!) \;\approx\; n \ln n - n + \frac{1}{2} \ln(2 \pi n).
        \]
        We can write this in terms of base-2 logarithms:
        \[
          \log_2(n!) \;=\; \frac{\ln(n!)}{\ln 2} \;\approx\; \frac{1}{\ln 2} \left[ n \ln n - n + \frac{1}{2} \ln(2\pi n) \right].
        \]
        Let us isolate the dominant terms as $n \to \infty$:
        * The term $n \ln n$ grows much faster than $n$ and $\ln n$.
        * Thus, the asymptotic behavior is dominated by the first term:
          \[
            \lim_{n \to \infty} \frac{\log_2(n!)}{n \log_2 n} \;=\; \lim_{n \to \infty} \frac{\frac{1}{\ln 2} (n \ln n - n + o(n))}{\frac{1}{\ln 2} (n \ln n)} \;=\; 1.
          \]
        By definition of $\Theta$-notation, this limit proves that:
        \[
          \log_2(n!) \;=\; \Theta(n \log n).
        \]
        This completes the proof. Since the decision tree height $h \ge \log_2(n!)$, the worst-case number of comparisons for any comparison-based sort must be $\Omega(n \log n)$.
\end{enumerate}
\end{problembox}

\newpage

% ─── Section: Problems ───────────────────────────────────────────────────────
\section{Part 6 --- Problems}

\begin{problembox}
\begin{enumerate}[label=\textbf{P\arabic*.}]
  \item \textbf{Solving Recurrences.}
  
        Solve the following recurrence relations using the Master Theorem or explain why it does not apply. If it does not apply, use the recurrence tree method to find a tight asymptotic bound:
        \begin{enumerate}[label=(\alph*)]
          \item $T(n) = 3 T(n/3) + n$.
          \item $T(n) = 8 T(n/2) + n^3$.
          \item $T(n) = 2 T(n/2) + n \log n$.
        \end{enumerate}

  \item \textbf{Comparison-Free Sorting (Counting Sort).}
  
        Comparison-based sorting has a lower bound of $\Omega(n \log n)$.
        However, the **Counting Sort** algorithm sorts $n$ integers in the range $[0, k]$ in $O(n + k)$ time.
        \begin{enumerate}[label=(\alph*)]
          \item Explain the mechanism of Counting Sort. (How does it use an auxiliary array to determine elements positions without comparing them?)
          \item If $k = O(n)$ (the range of integers is proportional to the list size), Counting Sort runs in linear $O(n)$ time.
                Explain why this does not violate the $\Omega(n \log n)$ lower bound.
                *(What assumption about the input data distinguishes Counting Sort from comparison-based decision trees?)*
        \end{enumerate}
\end{enumerate}
\end{problembox}

\newpage

% ─── Section: Deep Problems ──────────────────────────────────────────────────
\section{Part 7 --- Deep Problems}

\begin{deepbox}
\begin{enumerate}[label=\textbf{D\arabic*.}]
  \item \textbf{Average-Case Complexity of Quicksort.}
  
        Let $T(n)$ be the average number of comparisons executed by Quicksort to sort $n$ distinct elements.
        Assuming the pivot is chosen uniformly at random, the pivot splits the array into sizes $i$ and $n-1-i$ with equal probability $1/n$ for any $i \in [0, n-1]$.
        This yields the recurrence:
        \[
          T(n) \;=\; (n-1) + \frac{2}{n} \sum_{i=0}^{n-1} T(i)
        \]
        where $T(0) = T(1) = 0$.
        
        \begin{enumerate}[label=(\alph*)]
          \item Multiply the equation by $n$ to get:
                \[
                  n \cdot T(n) \;=\; n(n-1) + 2 \sum_{i=0}^{n-1} T(i).
                \]
                Write down the corresponding equation for $(n-1) \cdot T(n-1)$, and subtract it to prove the relation:
                \[
                  n \cdot T(n) \;=\; (n+1) T(n-1) + 2(n-1).
                \]
          \item Divide this by $n(n+1)$ to derive the telescoping sum:
                \[
                  \frac{T(n)}{n+1} \;=\; \frac{T(n-1)}{n} + \frac{2(n-1)}{n(n+1)} \;\le\; \frac{T(n-1)}{n} + \frac{2}{n}.
                \]
          \item Expand this telescoping sum to prove that:
                \[
                  T(n) \;\le\; 2 (n+1) \sum_{i=1}^n \frac{1}{i}.
                \]
                Using the integral approximation of the harmonic series ($\sum_{i=1}^n 1/i \approx \ln n$), prove that the average-case complexity of quicksort is $T(n) \le 2 n \ln n \approx 1.39 n \log_2 n \in \Theta(n \log n)$.
        \end{enumerate}

  \item \textbf{Linear-Time Selection (Median of Medians).}
  
        The selection problem is to find the $i$-th smallest element in an unsorted list of $n$ elements (e.g. finding the median).
        The **Median of Medians** algorithm guarantees linear $O(n)$ running time by recursively finding a good pivot:
        1. Divide the $n$ elements into $\lceil n/5 \rceil$ groups of 5.
        2. Find the median of each group.
        3. Recursively call the selection algorithm to find the median $x$ of these $\lceil n/5 \rceil$ medians.
        4. Partition the array around $x$ and recurse on the appropriate side.
        
        \begin{enumerate}[label=(\alph*)]
          \item Show that at least half of the $\lceil n/5 \rceil$ group medians are less than or equal to $x$.
                Conclude that at least $3n/10 - 6$ elements are less than or equal to $x$.
          \item Explain why this guarantees that the recursive call on the partition has size at most $7n/10 + 6$.
          \item Show that the total running time satisfies the recurrence:
                \[
                  T(n) \;\le\; T(\lceil n/5 \rceil) + T(7n/10 + 6) + c \cdot n.
                \]
                Prove by induction that $T(n) = O(n)$, showing that we can find the median of an unsorted list in linear time without sorting it.
        \end{enumerate}
\end{enumerate}
\end{deepbox}

\newpage

% ─── Section: Reflection ─────────────────────────────────────────────────────
\section{Part 8 --- Reflection}

\begin{reflectionbox}
\textbf{Conceptual Reflection}

Consider the following questions:
\begin{enumerate}
  \item **On the Limit of Infinite Scale:** Asymptotic analysis measures performance as inputs grow to infinity ($n \to \infty$). In the physical world, computer inputs are always finite, restricted by memory and lifespan. Why is an infinite limit still the best way to evaluate finite algorithm behavior?
  \item **On the Power of Randomization:** Randomized Quicksort uses a random number generator to bypass its worst-case input configurations. How does introducing "meaningless" entropy (randomness) into a deterministic machine prevent worst-case calculation failures?
  \item **On the Value of Information:** We proved that sorting requires at least $\Omega(n \log n)$ comparisons because each comparison yields at most 1 bit of information (yes/no). How does this lower bound demonstrate a fundamental connection between algorithmic complexity and information theory?
\end{enumerate}
\end{reflectionbox}

\medskip
\noindent\textbf{Looking ahead.}
Module~5.2 introduces **Advanced Data Structures and Balanced Trees**. Having analyzed how to sort and search data, we will study how to store it dynamically. We will study hash tables (universal hashing), self-balancing binary search trees (Red-Black trees and AVL rotations), and priority queues.

\vspace{1.4cm}
\begin{center}
{\color{JourneyColor}\rule{0.5\linewidth}{0.8pt}}\\[6pt]
{\small\color{JourneyColor}\textit{End of Journey 5, Module 5.1 --- Algorithmic Analysis and Divide-and-Conquer}}\\[4pt]
{\small\color{JourneyColor}Continue to: \textbf{Module 5.2 --- Advanced Data Structures and Balanced Trees}}\\[6pt]
{\color{JourneyColor}\rule{0.5\linewidth}{0.8pt}}
\end{center}

\end{document}
